\documentclass[11pt]{article}
\usepackage{mathblatt}


\blattfuss{Wahrscheinlichkeit}{Lösungen}
\begin{document}
\noindent{\large\bfseries Wahrscheinlichkeit \textperiodcentered{} Lösungen}\par\medskip
\begin{pfloesung}{}
\lz{1.}{$75\,\%$}{$\tfrac{3}{4}$}{}
\end{pfloesung}
\begin{pfloesung}{}
\lz{2.}{3 Kugeln}{}{}
\end{pfloesung}
\begin{pfloesung}{}
\lz{3.}{$0{,}5$}{$\tfrac{4}{8}$}{}
\end{pfloesung}
\begin{pfloesung}{}
\lz{4.}{60 Lose}{}{}
\end{pfloesung}
\begin{pfloesung}{}
\lz{5.}{$\tfrac{1}{6}$}{}{}
\end{pfloesung}
\begin{pfloesung}{}
\lz{6.}{4 Paare: , , ,}{grau; gelb; grau; weiß; grün; gelb; grün; weiß}{}
\end{pfloesung}
\begin{pfloesung}{}
\lz{7.}{6 Becher: , , , , ,}{Vanille; Karamell; Vanille; Himbeere; Schoko; Karamell; Schoko; Himbeere; Nuss; Karamell; Nuss; Himbeere}{}
\end{pfloesung}
\begin{pfloesung}{}
\lz{8.}{$6$ Kombinationen}{$3 \cdot 2$}{}
\end{pfloesung}
\begin{pfloesung}{}
\lz{9.}{6 Paare: , , , , ,}{1; 4; 2; 4; 3; 4; 4; 4; 5; 4; 6; 4}{}
\end{pfloesung}
\begin{pfloesung}{}
\lz{10.}{9 Zahlen: 44, 45, 46, 54, 55, 56, 64, 65, 66}{}{}
\end{pfloesung}
\begin{pfloesung}{{\small\color{mbgrau}P10 2025 · 3a}}
\lz{11.}{11, 12, 13, 21, 22, 23, 31, 32, 33}{3 · 3 = 9 Zahlen}{}
\end{pfloesung}
\begin{pfloesung}{}
\lz{12.}{$\tfrac{1}{3}$}{6 Zahlen: 147, 174, 417, 471, 714, 741; gerade sind 174 und 714, Anteil $\tfrac{2}{6}$}{}
\end{pfloesung}
\begin{pfloesung}{}
\lz{13.}{9 Ergebnisse: , , , , , , , ,}{5; 5; 5; 7; 5; 9; 7; 5; 7; 7; 7; 9; 9; 5; 9; 7; 9; 9}{}
\end{pfloesung}
\begin{pfloesung}{}
\lz{14.}{3 Ergebnisse: , ,}{Kopf; Kopf; Kopf; Zahl; Zahl; Kopf}{}
\end{pfloesung}
\begin{pfloesung}{{\small\color{mbgrau}P10 2017 · 1f}}
\lz{15.}{4 (ZZ, ZW, WZ, WW)}{}{}
\end{pfloesung}
\begin{pfloesung}{{\small\color{mbgrau}P10 2020 · 6a}}
\lz{16.}{(1; 2), (2; 2), (3; 2), (4; 2), (5; 2), (6; 2); $\tfrac{1}{6}$ $\approx$ 16,7 \%}{6 von 36 Augenpaaren}{}
\end{pfloesung}
\begin{pfloesung}{}
\lz{17.}{$\tfrac{1}{2}$}{}{}
\end{pfloesung}
\begin{pfloesung}{}
\lz{18.}{möglich: 35 Lose , günstig: 5 Gewinne}{30 Nieten und 5 Gewinne}{}
\end{pfloesung}
\begin{pfloesung}{}
\lz{19.}{$\tfrac{1}{6}$}{}{}
\end{pfloesung}
\begin{pfloesung}{}
\lz{20.}{$\tfrac{1}{2}$}{$\tfrac{3}{6}$}{}
\end{pfloesung}
\begin{pfloesung}{}
\lz{21.}{$\tfrac{1}{3}$}{$\tfrac{2}{6}$}{}
\end{pfloesung}
\begin{pfloesung}{{\small\color{mbgrau}P10 2014 · 1b}}
\lz{22.}{$\tfrac{1}{5}$ = 20 \%}{}{}
\end{pfloesung}
\begin{pfloesung}{{\small\color{mbgrau}P10 2014 · 1d}}
\lz{23.}{$\tfrac{2}{3}$}{}{}
\end{pfloesung}
\begin{pfloesung}{}
\lz{24.}{Nein: Beide Augenzahlen haben die Wahrscheinlichkeit $\tfrac{1}{6}$}{}{}
\end{pfloesung}
\begin{pfloesung}{}
\lz{25.}{Nein: $\tfrac{3}{16}$, denn es sind nur noch 16 Muffins da, davon 3 mit Chili}{}{}
\end{pfloesung}
\begin{pfloesung}{}
\lz{26.}{mittlere Dose}{}{}
\end{pfloesung}
\begin{pfloesung}{{\small\color{mbgrau}P10 2015 · 1a}}
\lz{27.}{rechter Topf (3 weiße, 3 graue Kugeln)}{}{}
\end{pfloesung}
\begin{pfloesung}{}
\lz{28.}{$\tfrac{1}{3}$}{$\tfrac{3}{9}$}{}
\end{pfloesung}
\begin{pfloesung}{}
\lz{29.}{$18\,\%$}{$\tfrac{9}{50} = 0{,}18$}{}
\end{pfloesung}
\begin{pfloesung}{}
\lz{30.}{$\tfrac{1}{4}$}{$\tfrac{15}{45 + 15} = \tfrac{15}{60}$}{}
\end{pfloesung}
\begin{pfloesung}{{\small\color{mbgrau}P10 2016 · 1f}}
\lz{31.}{$\tfrac{1}{20}$ = 5 \%}{}{}
\end{pfloesung}
\begin{pfloesung}{{\small\color{mbgrau}P10 2019 · 6a}}
\lz{32.}{$\tfrac{3}{20}$ = 15 \%}{}{}
\end{pfloesung}
\begin{pfloesung}{{\small\color{mbgrau}P10 2024 · 5a}}
\lz{33.}{$\tfrac{1}{5}$ = 20 \%}{1 günstiger von 5 Zetteln}{}
\end{pfloesung}
\begin{pfloesung}{{\small\color{mbgrau}P10 2017 · 6a}}
\lz{34.}{$\tfrac{1}{10}$ = 10 \%}{1 von 10 Ziffern}{}
\end{pfloesung}
\begin{pfloesung}{}
\lz{35.}{$0{,}35$}{$\tfrac{35}{100} = \tfrac{7}{20}$}{}
\end{pfloesung}
\begin{pfloesung}{{\small\color{mbgrau}P10 2016 · 5c}}
\lz{36.}{900 $-$ 101 + 1 = 800 Lose, 1 Hauptgewinn $\Rightarrow$ $\tfrac{1}{800}$}{800 Lose}{}
\end{pfloesung}
\begin{pfloesung}{{\small\color{mbgrau}P10 2014 · 6a}}
\lz{37.}{$\tfrac{3}{8}$; 37,5 \%}{3 von 8 Feldern}{}
\end{pfloesung}
\begin{pfloesung}{{\small\color{mbgrau}P10 2016 · 5d}}
\lz{38.}{ja, 7 Trostpreis-Lose unter 80 Losen mit Endziffer 6 $\Rightarrow$ $\tfrac{7}{80}$}{80 Lose mit Endziffer 6; 8 mit Endung 26, ohne 326 $\Rightarrow$ 7}{}
\end{pfloesung}
\begin{pfloesung}{{\small\color{mbgrau}P10 2018 · 7b}}
\lz{39.}{nein, P = $\tfrac{2}{14}$ = $\tfrac{1}{7}$ $\approx$ 14,3 \%, nicht $\tfrac{2}{16}$}{noch 14 Pfannkuchen, 2 mit Senf}{}
\end{pfloesung}
\begin{pfloesung}{}
\lz{40.}{8 Kekse, davon 3 mit Schoko}{}{}
\end{pfloesung}
\begin{pfloesung}{}
\lz{41.}{mit Zurücklegen}{}{}
\end{pfloesung}
\begin{pfloesung}{}
\lz{42.}{auf beiden Stufen rot $\tfrac{1}{3}$, blau $\tfrac{2}{3}$}{}{}
\end{pfloesung}
\begin{pfloesung}{}
\lz{43.}{$\tfrac{3}{28}$}{$\tfrac{3}{8} \cdot \tfrac{2}{7} = \tfrac{6}{56}$}{}
\end{pfloesung}
\begin{pfloesung}{}
\lz{44.}{das Glücksrad ändert sich nicht}{an jedem Punkt grün $\tfrac{3}{5}$, gelb $\tfrac{2}{5}$}{}
\end{pfloesung}
\begin{pfloesung}{}
\lz{45.}{nach rot: rot $\tfrac{3}{6}$}{erste Stufe grün $\tfrac{3}{7}$}{}
\end{pfloesung}
\begin{pfloesung}{}
\lz{46.}{erste Stufe ohne Nuss $\tfrac{10}{15}$}{}{}
\end{pfloesung}
\begin{pfloesung}{{\small\color{mbgrau}P10 2024 · 5b}}
\lz{47.}{1. Woche: kein Joker 4/5, Joker $\tfrac{1}{5}$; 2. Woche an beiden Ästen: kein Joker 4/5, Joker $\tfrac{1}{5}$; P(Joker, Joker) = $\tfrac{1}{5}$ · $\tfrac{1}{5}$ = $\tfrac{1}{25}$}{mit Zurücklegen: jede Woche $\tfrac{1}{5}$}{}
\end{pfloesung}
\begin{pfloesung}{}
\lz{48.}{z. B. ein Glücksrad mit fünf gleich großen Feldern, zwei rot und drei blau, zweimal gedreht}{}{}
\end{pfloesung}
\begin{pfloesung}{{\small\color{mbgrau}P10 2014 · 6b}}
\lz{49.}{1. Stufe A 1/2, B $\tfrac{3}{8}$; 2. Stufe je Ast A 1/2, B 3/8, C $\tfrac{1}{8}$}{}{}
\end{pfloesung}
\begin{pfloesung}{{\small\color{mbgrau}P10 2019 · 6b}}
\lz{50.}{$\tfrac{16}{20}$; $\tfrac{3}{18}$; Aussage 1 falsch, 2 falsch, 3 wahr}{fetter Pfad = dreimal rot; $\tfrac{4}{20}$ · $\tfrac{3}{19}$ · $\tfrac{2}{18}$ = $\tfrac{1}{285}$ $\approx$ 0,35 \%}{}
\end{pfloesung}
\begin{pfloesung}{{\small\color{mbgrau}P10 2020 · 6c}}
\lz{51.}{Äste 6: 1/6, keine 6: $\tfrac{5}{6}$; $\tfrac{2}{27}$ $\approx$ 7,4 \%}{3 · ($\tfrac{1}{6}$)² · $\tfrac{5}{6}$ = $\tfrac{15}{216}$; ($\tfrac{1}{6}$)³ = $\tfrac{1}{216}$}{}
\end{pfloesung}
\begin{pfloesung}{{\small\color{mbgrau}P10 2015 · 7d}}
\lz{52.}{10/11, 1/10, 1/9, $\tfrac{8}{9}$; $\tfrac{3}{11}$ $\approx$ 0,27}{1 $-$ $\tfrac{10}{11}$ · $\tfrac{9}{10}$ · $\tfrac{8}{9}$; jeder Pfad mit R hat $\tfrac{1}{11}$}{}
\end{pfloesung}
\begin{pfloesung}{}
\lz{53.}{Gegenereignis: keine Sechs in beiden Würfen}{}{}
\end{pfloesung}
\begin{pfloesung}{}
\lz{54.}{2 Pfade: ,}{rot; blau; blau; rot}{}
\end{pfloesung}
\begin{pfloesung}{}
\lz{55.}{$\tfrac{125}{216}$}{$\tfrac{5}{6} \cdot \tfrac{5}{6} \cdot \tfrac{5}{6}$}{}
\end{pfloesung}
\begin{pfloesung}{}
\lz{56.}{$\tfrac{9}{64}$}{$\tfrac{3}{8} \cdot \tfrac{3}{8}$}{}
\end{pfloesung}
\begin{pfloesung}{}
\lz{57.}{$\tfrac{7}{16}$}{$\tfrac{3}{4} \cdot \tfrac{2}{4} + \tfrac{1}{4} \cdot \tfrac{1}{4} = \tfrac{6}{16} + \tfrac{1}{16}$}{}
\end{pfloesung}
\begin{pfloesung}{}
\lz{58.}{$\tfrac{15}{32}$}{$\tfrac{25}{64} + \tfrac{4}{64} + \tfrac{1}{64} = \tfrac{30}{64}$}{}
\end{pfloesung}
\begin{pfloesung}{}
\lz{59.}{Nein: Zwei gleiche Augenzahlen haben $\tfrac{6}{36} = \tfrac{1}{6}$, eine 1 im ersten Wurf hat auch $\tfrac{1}{6}$}{}{}
\end{pfloesung}
\begin{pfloesung}{{\small\color{mbgrau}P10 2025 · 3b}}
\lz{60.}{P(33) = $\tfrac{1}{4}$ · $\tfrac{1}{4}$ = $\tfrac{1}{16}$; falsch, P(22) = $\tfrac{1}{4}$ · $\tfrac{2}{4}$ = $\tfrac{1}{8}$ $\neq$ $\tfrac{1}{16}$}{P(3 links) = $\tfrac{1}{4}$; P(2 rechts) = $\tfrac{2}{4}$}{}
\end{pfloesung}
\begin{pfloesung}{}
\lz{61.}{$0{,}64$}{$0{,}8 \cdot 0{,}8$}{}
\end{pfloesung}
\begin{pfloesung}{{\small\color{mbgrau}P10 2020 · 6b}}
\lz{62.}{Max hat nicht Recht; beide Wahrscheinlichkeiten $\tfrac{1}{6}$}{P(gleich) = $\tfrac{6}{36}$ = $\tfrac{1}{6}$; P(zuerst 6) = $\tfrac{1}{6}$}{}
\end{pfloesung}
\begin{pfloesung}{{\small\color{mbgrau}P10 2025 · 3c}}
\lz{63.}{$\tfrac{5}{16}$ = 31,25 \%}{P(12) = $\tfrac{4}{16}$; P(21) = $\tfrac{1}{16}$}{}
\end{pfloesung}
\begin{pfloesung}{{\small\color{mbgrau}P10 2014 · 6c}}
\lz{64.}{$\tfrac{13}{32}$ $\approx$ 0,41}{$\tfrac{1}{4}$ + $\tfrac{9}{64}$ + $\tfrac{1}{64}$ = $\tfrac{26}{64}$}{}
\end{pfloesung}
\begin{pfloesung}{}
\lz{65.}{eine Stufe}{}{}
\end{pfloesung}
\begin{pfloesung}{}
\lz{66.}{drei Stufen}{}{}
\end{pfloesung}
\begin{pfloesung}{}
\lz{67.}{$\tfrac{2}{5}$}{$\tfrac{3}{5} \cdot \tfrac{2}{4} + \tfrac{2}{5} \cdot \tfrac{1}{4} = \tfrac{6}{20} + \tfrac{2}{20} = \tfrac{8}{20}$}{}
\end{pfloesung}
\begin{pfloesung}{}
\lz{68.}{$\tfrac{1}{7}$}{$\tfrac{6}{7} \cdot \tfrac{1}{6}$}{}
\end{pfloesung}
\begin{pfloesung}{}
\lz{69.}{$\tfrac{1}{21}$}{$\tfrac{4}{9} \cdot \tfrac{3}{8} \cdot \tfrac{2}{7} = \tfrac{24}{504}$}{}
\end{pfloesung}
\begin{pfloesung}{}
\lz{70.}{$\tfrac{1}{3}$}{$\tfrac{1}{9} + \tfrac{8}{9} \cdot \tfrac{1}{8} + \tfrac{8}{9} \cdot \tfrac{7}{8} \cdot \tfrac{1}{7} = \tfrac{3}{9}$}{}
\end{pfloesung}
\begin{pfloesung}{}
\lz{71.}{zehnmal so wahrscheinlich, nicht doppelt}{Drei rote: $\tfrac{5}{15} \cdot \tfrac{4}{14} \cdot \tfrac{3}{13} = \tfrac{60}{2\,730}$, drei grüne: $\tfrac{3}{15} \cdot \tfrac{2}{14} \cdot \tfrac{1}{13} = \tfrac{6}{2\,730}$}{}
\end{pfloesung}
\begin{pfloesung}{}
\lz{72.}{$0{,}7$}{$1 - \tfrac{3}{5} \cdot \tfrac{2}{4} = 1 - \tfrac{6}{20}$}{}
\end{pfloesung}
\begin{pfloesung}{{\small\color{mbgrau}P10 2024 · 5c}}
\lz{73.}{$\tfrac{1}{5}$}{$\tfrac{4}{5}$ · $\tfrac{1}{4}$ = $\tfrac{4}{20}$}{}
\end{pfloesung}
\begin{pfloesung}{{\small\color{mbgrau}P10 2019 · 6c}}
\lz{74.}{$\tfrac{1}{57}$ $\approx$ 1,75 \% gegen $\tfrac{1}{1140}$ $\approx$ 0,09 \% $\Rightarrow$ 20-mal so hoch, Behauptung falsch}{$\tfrac{6}{20}$ · $\tfrac{5}{19}$ · $\tfrac{4}{18}$ = $\tfrac{1}{57}$; $\tfrac{3}{20}$ · $\tfrac{2}{19}$ · $\tfrac{1}{18}$ = $\tfrac{1}{1140}$}{}
\end{pfloesung}
\begin{pfloesung}{{\small\color{mbgrau}P10 2018 · 7c}}
\lz{75.}{(S; M), (M; S), (S; S); $\tfrac{91}{120}$ $\approx$ 0,758; beide Pfannkuchen mit derselben Füllung}{$\tfrac{14}{16}$ · $\tfrac{13}{15}$ = $\tfrac{182}{240}$}{}
\end{pfloesung}
\begin{pfloesung}{}
\lz{76.}{$\tfrac{1}{2}$}{$\tfrac{3}{6}$}{}
\end{pfloesung}
\begin{pfloesung}{}
\lz{77.}{$\tfrac{11}{36}$}{$1 - \tfrac{5}{6} \cdot \tfrac{5}{6} = 1 - \tfrac{25}{36}$}{}
\end{pfloesung}
\begin{pfloesung}{}
\lz{78.}{$\tfrac{7}{12}$}{$1 - \tfrac{3}{6} \cdot \tfrac{5}{6} = 1 - \tfrac{15}{36} = \tfrac{21}{36}$}{}
\end{pfloesung}
\begin{pfloesung}{}
\lz{79.}{$3$ Felder}{$0{,}75 \cdot 4$}{}
\end{pfloesung}
\begin{pfloesung}{}
\lz{80.}{3 rote, 2 blaue, 1 gelbes Feld}{}{}
\end{pfloesung}
\begin{pfloesung}{}
\lz{81.}{$\tfrac{3}{8}$}{z. B. links 5, 5, 5, 6 und rechts 7, 7, 8, 8: $\tfrac{3}{4} \cdot \tfrac{2}{4} = \tfrac{6}{16}$}{}
\end{pfloesung}
\begin{pfloesung}{{\small\color{mbgrau}P10 2018 · 1j}}
\lz{82.}{2 rote Felder}{}{}
\end{pfloesung}
\begin{pfloesung}{{\small\color{mbgrau}P10 2025 · 3d}}
\lz{83.}{z. B. links 1, 1, 2, 3 und rechts 3, 3, 1, 2: P(13) = $\tfrac{2}{4}$ · $\tfrac{2}{4}$ = $\tfrac{1}{4}$ = 25 \%}{$\tfrac{1}{4}$ = $\tfrac{2}{4}$ · $\tfrac{2}{4}$ (oder 1 · $\tfrac{1}{4}$)}{}
\end{pfloesung}
\begin{pfloesung}{{\small\color{mbgrau}P10 2019 · 1g}}
\lz{84.}{2 weiße Kugeln}{6 Kugeln insgesamt}{}
\end{pfloesung}
\begin{pfloesung}{\textbf{Prüfstein} \textperiodcentered{} P10 2026 FOR · 6}
\lz{a)}{$\tfrac{2}{6}$ = $\tfrac{1}{3}$}{2 günstige von 6 Flächen}{}
\lz{b)}{A ungerade $\tfrac{3}{6}$; B nach A gerade: ungerade $\tfrac{5}{6}$; B nach A ungerade: gerade 1/6, ungerade $\tfrac{5}{6}$; P(beide ungerade) = $\tfrac{3}{6}$ · $\tfrac{5}{6}$ = $\tfrac{15}{36}$ = $\tfrac{5}{12}$}{P(B gerade) = $\tfrac{1}{6}$; P(B ungerade) = $\tfrac{5}{6}$}{}
\lz{c)}{$\tfrac{33}{36}$ = $\tfrac{11}{12}$}{P(gerade, gerade) = $\tfrac{3}{6}$ · $\tfrac{1}{6}$ = $\tfrac{3}{36}$}{}
\lz{d)}{z. B. B neu: 1, 1, 1, 3, 3, 3 (die 2 durch 1 ersetzt), P(1, 1) = $\tfrac{2}{6}$ · $\tfrac{3}{6}$ = $\tfrac{6}{36}$ = $\tfrac{1}{6}$}{Summe 2 nur bei 1 + 1; P(A = 1) = $\tfrac{2}{6}$ $\Rightarrow$ P(B = 1) = $\tfrac{3}{6}$}{}
\end{pfloesung}
\end{document}